% answers_short.tex —— 1990 年·广东卷·理科 · 简短答案速查
% 由 tools/make_short_answers.py 自 answers.tex 的 \gksolhead/\gkanswer 提取，
% 并与 1990/1990_数学真题_广东理_答案要点.md 逐题比对后生成，勿手改（改源 MD 后重跑）。

\begin{gknotice}
\textbf{使用说明}：本节只给最终答案便于快速核对，完整推导见解析版。本卷包含来源印载答案与项目整理的推导，尚未经过异人盲核。第12题题图缺失，B仅为来源印载答案，数学结论待核。\textbf{非官方标准答案}。
\end{gknotice}

\gkpart{选择题}
\begin{center}\begin{tabular}{cccccccccccccc}
\toprule
题号 & 1 & 2 & 3 & 4 & 5 & 6 & 7 & 8 & 9 & 10 & 11 & 12 & 13 \\
\midrule
答案 & D & A & C & C & C & D & C & B & B & A & A & B & B \\
\bottomrule
\end{tabular}\end{center}\medskip
\begin{center}\begin{tabular}{ccccccccccccc}
\toprule
题号 & 14 & 15 & 16 & 17 & 18 & 19 & 20 & 21 & 22 & 23 & 24 & 25 \\
\midrule
答案 & C & A & D & A & C & B & B & D & B & C & D & D \\
\bottomrule
\end{tabular}\end{center}

\gkpart{解答题}
\begin{itemize}[leftmargin=2.6em]
  \item[\textbf{一.}] (1) $x=\pm\dfrac{16}{5}$；(2) $-10$；(3) $2$；(4) $\sqrt2$；(5) $7:5$
  \item[\textbf{二.}] (1) 见解析版；(2) $60^\circ$
  \item[\textbf{三.}] (1) $\dfrac{y^2}{4}+\dfrac{x^2}{3}=1$；(2) $\operatorname{tg}\angle F_1PF_2=\dfrac43$
  \item[\textbf{四.}] (1) $a_1=\dfrac\pi4$，$d=\dfrac\pi3$；(2) $b_n=\dfrac{\sqrt2}{8}(-1)^{n-1}$
  \item[\textbf{五.}] (1) 见解析版；(2) $a>-\dfrac34$
\end{itemize}
